% Generated draft: requires mathematical and source review.
\documentclass[11pt]{article}
\usepackage[margin=25mm]{geometry}
\usepackage{fontspec}
\setmainfont{Latin Modern Roman}
\usepackage{amsmath,amssymb,mathtools,unicode-math}
\setmathfont{Latin Modern Math}
\usepackage{enumitem,needspace,xurl,hyperref}
\setlength{\parindent}{0pt}
\setlength{\parskip}{5pt}
\setlength{\emergencystretch}{3em}
\newcommand{\N}{\mathbb{N}}
\newcommand{\Z}{\mathbb{Z}}
\newcommand{\Q}{\mathbb{Q}}
\newcommand{\R}{\mathbb{R}}
\newcommand{\C}{\mathbb{C}}
\newcommand{\F}{\mathbb{F}}
\newcommand{\A}{\mathbb{A}}
\newcommand{\PP}{\mathbb{P}}
\newcommand{\E}{\mathbb{E}}
\newcommand{\eps}{\varepsilon}
\newcommand{\dd}{\,\mathrm{d}}
\begin{document}
\section*{UnsolvedMath ID: 801184}
\textbf{Source problem code:} catalogue entry
\subsection{Definitions and mathematical statement}
Let \(S_\infty=\bigcup_{n\geq 1}S_n\), where \(S_n\hookrightarrow S_{n+1}\) fixes \(n+1\); thus its elements are permutations of the positive integers with finite support. Write \(s_i=(i,i+1)\) and let \(\ell(w)\) be the number of inversions of \(w\). On \(\Z[x_1,x_2,\ldots]\), define the divided difference operator
\[
\partial_i f=\frac{f-s_i f}{x_i-x_{i+1}},
\]
where \(s_i f\) exchanges the variables \(x_i\) and \(x_{i+1}\). For \(w\in S_n\), the Schubert polynomial \(\mathfrak S_w\) is characterized by
\[
\mathfrak S_{w_0^{(n)}}=x_1^{n-1}x_2^{n-2}\cdots x_{n-1},\qquad
\mathfrak S_w=\partial_i\mathfrak S_{ws_i}
\quad\text{if }\ell(ws_i)=\ell(w)+1,
\]
where \(w_0^{(n)}\) is the longest permutation in \(S_n\). These definitions agree under the inclusions fixing trailing letters.

The Schubert polynomials form a \(\Z\)-basis of \(\Z[x_1,x_2,\ldots]\). Define the integers \(c_{uv}^{w}\), for \(u,v,w\in S_\infty\), by the unique finite expansion
\[
\mathfrak S_u\mathfrak S_v
=\sum_{w\in S_\infty}c_{uv}^{w}\mathfrak S_w.
\]
In particular, \(c_{uv}^{w}=0\) unless \(\ell(w)=\ell(u)+\ell(v)\).

\textbf{Problem (Stanley, Problem 11, as supplied).} Give a uniform positive combinatorial interpretation of these structure constants: specify finite sets \(\mathcal C(u,v;w)\), defined by explicit combinatorial conditions on \(u,v,w\), and prove that
\[
\forall u,v,w\in S_\infty,\qquad
c_{uv}^{w}=\lvert\mathcal C(u,v;w)\rvert.
\]
The construction should supply a combinatorial proof of nonnegativity, rather than define the sets using the already known coefficients. Here ``positive'' means a counting interpretation yielding nonnegative integers, including zero.

The supplied record states that nonnegativity is known through intersection theory on flag varieties \(\operatorname{GL}(n,\C)/B\), where \(B\) is the subgroup of invertible upper triangular matrices. It labels the general problem open and reports known special cases described as Grassmannian, two-step flag, and separated descents. These are supplied status qualifications, not an independently verified literature assessment. The record does not specify the precise permutation hypotheses or references for those special cases; consequently, they are retained as contextual qualifications rather than additional formal assertions.
\subsection{Short English statement}
Find an explicit combinatorial counting rule for every coefficient in the product of two Schubert polynomials. Such a rule should explain combinatorially why all coefficients are nonnegative. The supplied record attributes the problem to Stanley in 1999, labels it open in general, and reports special cases whose exact scope is unspecified.
\subsection{Sources}
\begin{itemize}
\item Local UnsolvedMath registry, supplied complete source record for permanent ID \(801184\), source problem code \texttt{catalogue entry}, in the AMR Open Problem Lists collection. The record attributes the problem to Stanley (1999) and labels its status \texttt{open}.
\item Richard Stanley, \emph{Positivity Problems and Conjectures in Algebraic Combinatorics} (1999), Problem 11, as cited in the supplied record: \url{https://math.mit.edu/~rstan/papers/problems.pdf}.
\item Supplied provenance records access on \texttt{2026-10-16} and extraction method \texttt{pdf+html\_verified}. The associated supplied curation reports known Grassmannian, two-step flag, and separated-descents cases and an open general case. These provenance and status assertions are supplied dated curation, not independently verified here; no separate \texttt{literature\_checked\_at} field or special-case literature references were supplied.
\end{itemize}
\end{document}
